% !TEX program = xelatex
\documentclass[aspectratio=169]{beamer}
\usetheme{UFS}
\usepackage[utf8]{inputenc}
\usepackage[portuguese]{babel}
\usepackage{amsmath,amssymb}
\usepackage{graphicx}
\usepackage{booktabs}
\usepackage{listings}
\usepackage{xcolor}
\usepackage{tikz}
\usetikzlibrary{shapes.geometric, arrows.meta, positioning}

% ── VS Code Light+ palette ──
\definecolor{bgcolor}{HTML}{FFFFFF}
\definecolor{linecolor}{HTML}{DDDDDD}
\definecolor{numcolor}{HTML}{999999}
\definecolor{kwcolor}{HTML}{0000FF}
\definecolor{strcolor}{HTML}{A31515}
\definecolor{cmtcolor}{HTML}{008000}
\definecolor{fncolor}{HTML}{795E26}
\definecolor{typecolor}{HTML}{267F99}

\lstdefinestyle{modernstyle}{
  language=Python,
  basicstyle=\ttfamily\scriptsize,
  backgroundcolor=\color{bgcolor},
  rulecolor=\color{linecolor},
  numberstyle=\tiny\color{numcolor},
  keywordstyle=\color{kwcolor}\bfseries,
  stringstyle=\color{strcolor},
  commentstyle=\color{cmtcolor}\itshape,
  emphstyle=\color{fncolor},
  emphstyle={[2]\color{typecolor}},
  frame=single,
  numbers=left,
  numbersep=6pt,
  tabsize=4,
  showstringspaces=false,
  breaklines=true,
  columns=fullflexible,
  literate={á}{{\'a}}1 {é}{{\'e}}1 {í}{{\'i}}1 {ó}{{\'o}}1 {ú}{{\'u}}1
           {ã}{{\~a}}1 {õ}{{\~o}}1 {ç}{{\c{c}}}1 {Ç}{{\c{C}}}1
           {à}{{\`a}}1 {ê}{{\^e}}1,
  emph={from_csv_line,from_dict,calcular_taxa_servico,calcular_desconto,
        criar,registrar,adicionar,total,processar,finalizar,
        aplicar,gerar_relatorio,salvar,enviar_email},
  emph={[2]Lanche,Pedido,Hamburguer,XTudo,XSalada,Bebida,
        DescontoPercentual,DescontoFixo,SemDesconto,LancheFactory,
        Cozinha,Caixa,RelatorioHTML,Armazenamento,Enviador,
        PedidoRuim,Retangulo,Quadrado,ABC,abstractmethod,
        dataclass,field},
}
\lstset{style=modernstyle}

\title{COMP0496 -- Programação A \\ Programação Orientada a Objetos --- Parte 4}
\subtitle{Tópicos Avançados e Boas Práticas}
\author{Prof.\ Dr.\ André Yoshiaki Kashiwabara}
\institute{DCOMP -- Universidade Federal de Sergipe\\\texttt{andreyoshiaki@dcomp.ufs.br}}
\date{}

\begin{document}

\begin{frame}
\titlepage
\end{frame}

\begin{frame}{Roteiro}
\tableofcontents
\end{frame}

% ================================================================
\section{@dataclass --- Eliminando Boilerplate}
% ================================================================

\begin{frame}[fragile]{O Problema: Boilerplate Repetitivo}
Uma classe simples com \texttt{\_\_init\_\_}, \texttt{\_\_repr\_\_} e \texttt{\_\_eq\_\_}:

\begin{lstlisting}
class Lanche:
    def __init__(self, nome, preco, ingredientes):
        self.nome = nome
        self.preco = preco
        self.ingredientes = ingredientes

    def __repr__(self):
        return (f"Lanche(nome={self.nome!r}, "
                f"preco={self.preco!r}, "
                f"ingredientes={self.ingredientes!r})")

    def __eq__(self, other):
        if not isinstance(other, Lanche):
            return NotImplemented
        return (self.nome == other.nome
                and self.preco == other.preco
                and self.ingredientes == other.ingredientes)
\end{lstlisting}

\begin{alertblock}{Problema}
$\sim$20 linhas para guardar 3 campos. E se a classe tiver 8 campos?
\end{alertblock}
\end{frame}

\begin{frame}[fragile]{@dataclass: A Solução}
\begin{lstlisting}
from dataclasses import dataclass

@dataclass
class Lanche:
    nome: str
    preco: float
    ingredientes: list[str]
\end{lstlisting}

O decorador gera automaticamente:
\begin{itemize}
  \item \texttt{\_\_init\_\_} --- com todos os campos como parâmetros
  \item \texttt{\_\_repr\_\_} --- representação legível
  \item \texttt{\_\_eq\_\_} --- comparação campo a campo
\end{itemize}

\begin{exampleblock}{Resultado}
3 linhas substituem $\sim$20. O código expressa \textbf{o quê}, não \textbf{o como}.
\end{exampleblock}
\end{frame}

\begin{frame}[fragile]{Parâmetros Úteis do @dataclass}
\begin{lstlisting}
@dataclass(frozen=True)
class ItemCardapio:
    nome: str
    preco: float
    categoria: str
\end{lstlisting}

\begin{table}
\centering
\begin{tabular}{lll}
\toprule
\textbf{Parâmetro} & \textbf{Efeito} & \textbf{Quando usar} \\
\midrule
\texttt{frozen=True}  & Imutável (sem \texttt{setattr}) & Itens de cardápio, configs \\
\texttt{order=True}   & Gera \texttt{<, <=, >, >=}     & Ordenar por campos \\
\texttt{eq=True}       & Gera \texttt{\_\_eq\_\_} (padrão) & Quase sempre \\
\texttt{slots=True}   & Usa \texttt{\_\_slots\_\_} (3.10+) & Performance \\
\bottomrule
\end{tabular}
\end{table}
\end{frame}

\begin{frame}[fragile]{Defaults Mutáveis: field(default\_factory)}
\begin{lstlisting}
from dataclasses import dataclass, field

@dataclass
class Pedido:
    cliente: str
    itens: list[str] = field(default_factory=list)
    observacoes: dict[str, str] = field(default_factory=dict)
    numero: int = field(default=0, repr=False)
\end{lstlisting}

\begin{alertblock}{Nunca faça isso}
\texttt{itens: list[str] = []} --- o mesmo objeto é compartilhado entre instâncias!
\end{alertblock}

\texttt{field(default\_factory=list)} cria uma \textbf{nova lista} para cada instância.
\end{frame}

\begin{frame}[fragile]{@dataclass com Métodos Normais}
\begin{lstlisting}
@dataclass
class Pedido:
    cliente: str
    itens: list[str] = field(default_factory=list)

    def adicionar(self, item: str):
        self.itens.append(item)

    def total(self, cardapio: dict[str, float]) -> float:
        return sum(cardapio[i] for i in self.itens)

    def resumo(self) -> str:
        return f"{self.cliente}: {len(self.itens)} itens"
\end{lstlisting}

\begin{exampleblock}{Ponto-chave}
\texttt{@dataclass} gera o boilerplate; você adiciona a lógica de negócio.
\end{exampleblock}
\end{frame}

\begin{frame}{Quando Usar / Quando NÃO Usar @dataclass}
\begin{columns}[T]
\column{0.48\textwidth}
\textbf{Use quando:}
\begin{itemize}
  \item Classe é principalmente ``dados com métodos''
  \item Precisa de \texttt{\_\_eq\_\_} e \texttt{\_\_repr\_\_}
  \item Muitos campos (reduz boilerplate)
  \item Quer imutabilidade (\texttt{frozen})
\end{itemize}

\column{0.48\textwidth}
\textbf{NÃO use quando:}
\begin{itemize}
  \item \texttt{\_\_init\_\_} tem lógica complexa (validação pesada)
  \item Herança profunda ($>$2 níveis)
  \item Classe é majoritariamente comportamento
  \item Compatibilidade com Python $<$3.7
\end{itemize}
\end{columns}
\end{frame}

% ================================================================
\section{@classmethod e @staticmethod}
% ================================================================

\begin{frame}{Três Tipos de Método}
\begin{table}
\centering
\begin{tabular}{llll}
\toprule
\textbf{Tipo} & \textbf{Decorador} & \textbf{1º Parâmetro} & \textbf{Acessa} \\
\midrule
Instância   & (nenhum)          & \texttt{self} & instância + classe \\
Classe      & \texttt{@classmethod}  & \texttt{cls}  & apenas classe \\
Estático    & \texttt{@staticmethod} & (nenhum)      & nada da classe \\
\bottomrule
\end{tabular}
\end{table}

\begin{exampleblock}{Regra prática}
\textbf{Precisa de \texttt{self}?} $\to$ método de instância.
\textbf{Precisa de \texttt{cls}?} $\to$ classmethod.
\textbf{Não precisa de nenhum?} $\to$ staticmethod.
\end{exampleblock}
\end{frame}

\begin{frame}[fragile]{@classmethod como Construtor Alternativo}
\begin{lstlisting}
@dataclass
class Lanche:
    nome: str
    preco: float
    categoria: str

    @classmethod
    def from_csv_line(cls, linha: str) -> "Lanche":
        nome, preco, cat = linha.strip().split(",")
        return cls(nome, float(preco), cat)

    @classmethod
    def from_dict(cls, d: dict) -> "Lanche":
        return cls(**d)
\end{lstlisting}

\begin{lstlisting}
x = Lanche.from_csv_line("X-Tudo,25.90,hamburguer")
y = Lanche.from_dict({"nome": "Suco", "preco": 8.0,
                       "categoria": "bebida"})
\end{lstlisting}
\end{frame}

\begin{frame}[fragile]{Por Que Não Uma Função Comum?}
\begin{columns}[T]
\column{0.48\textwidth}
\textbf{Função livre:}
\begin{lstlisting}
def lanche_from_csv(linha):
    nome, preco, cat = \
        linha.split(",")
    return Lanche(
        nome, float(preco), cat
    )
\end{lstlisting}
\begin{itemize}
  \item Funciona, mas fica ``solta''
  \item Não herda: sempre cria \texttt{Lanche}
\end{itemize}

\column{0.48\textwidth}
\textbf{@classmethod:}
\begin{lstlisting}
@classmethod
def from_csv_line(cls, linha):
    nome, preco, cat = \
        linha.split(",")
    return cls(
        nome, float(preco), cat
    )
\end{lstlisting}
\begin{itemize}
  \item Agrupada na classe
  \item \texttt{cls} cria o subtipo correto
\end{itemize}
\end{columns}

\begin{alertblock}{Herança}
Se \texttt{XTudo(Lanche)}, então \texttt{XTudo.from\_csv\_line(...)} retorna um \texttt{XTudo}, não um \texttt{Lanche}.
\end{alertblock}
\end{frame}

\begin{frame}[fragile]{@staticmethod para Utilitários}
\begin{lstlisting}
class Lanche:
    # ...

    @staticmethod
    def calcular_taxa_servico(valor: float,
                               percentual: float = 0.10) -> float:
        """Taxa de serviço sobre o valor."""
        return round(valor * percentual, 2)
\end{lstlisting}

\begin{lstlisting}
taxa = Lanche.calcular_taxa_servico(50.00)  # 5.0
taxa = Lanche.calcular_taxa_servico(50.00, 0.15)  # 7.5
\end{lstlisting}

\begin{exampleblock}{Quando usar}
Lógica relacionada à classe, mas que não precisa de \texttt{self} nem de \texttt{cls}. Alternativa: função no módulo.
\end{exampleblock}
\end{frame}

\begin{frame}{Guia de Decisão}
\begin{center}
\begin{tikzpicture}[
  node distance=1.2cm,
  every node/.style={font=\small},
  decision/.style={diamond, draw, fill=blue!10, text width=3.2cm,
    align=center, inner sep=1pt, aspect=2},
  result/.style={rectangle, draw, fill=green!10, rounded corners,
    text width=3cm, align=center, minimum height=0.8cm},
  arr/.style={->, thick}
]
\node[decision] (q1) {Precisa acessar\\a instância?};
\node[result, right=2.5cm of q1] (r1) {Método de\\instância (\texttt{self})};
\node[decision, below=1.2cm of q1] (q2) {Precisa acessar\\a classe?};
\node[result, right=2.5cm of q2] (r2) {\texttt{@classmethod}\\(\texttt{cls})};
\node[result, below=1.2cm of q2] (r3) {\texttt{@staticmethod}};

\draw[arr] (q1) -- node[above]{Sim} (r1);
\draw[arr] (q1) -- node[left]{Não} (q2);
\draw[arr] (q2) -- node[above]{Sim} (r2);
\draw[arr] (q2) -- node[left]{Não} (r3);
\end{tikzpicture}
\end{center}
\end{frame}

% ================================================================
\section{Princípios SOLID (Introdução)}
% ================================================================

\begin{frame}{SOLID --- Visão Geral}
\begin{table}
\centering
\begin{tabular}{cll}
\toprule
\textbf{Letra} & \textbf{Princípio} & \textbf{Ideia Central} \\
\midrule
\textbf{S} & Single Responsibility  & Uma classe, uma razão para mudar \\
\textbf{O} & Open/Closed            & Aberta para extensão, fechada para modificação \\
\textbf{L} & Liskov Substitution    & Subtipos substituem o pai sem surpresas \\
\textbf{I} & Interface Segregation  & Interfaces pequenas e coesas \\
\textbf{D} & Dependency Inversion   & Dependa de abstrações, não de concretos \\
\bottomrule
\end{tabular}
\end{table}

Hoje veremos \textbf{S}, \textbf{O} e \textbf{L} com exemplos do Podrão.
\end{frame}

\begin{frame}[fragile]{SRP --- O Problema}
\begin{lstlisting}
class PedidoRuim:
    def __init__(self, cliente, itens):
        self.cliente = cliente
        self.itens = itens

    def calcular_total(self): ...       # lógica de negócio
    def gerar_relatorio_html(self): ... # apresentação
    def salvar_em_arquivo(self): ...    # persistência
    def enviar_email(self): ...         # comunicação
\end{lstlisting}

\begin{alertblock}{4 razões para mudar}
Mudança no cálculo, no formato HTML, no storage, ou no e-mail --- todas alteram a mesma classe.
\end{alertblock}
\end{frame}

\begin{frame}[fragile]{SRP --- A Solução: Separar Responsabilidades}
\begin{lstlisting}
class Pedido:
    def calcular_total(self): ...

class RelatorioHTML:
    def gerar(self, pedido: Pedido): ...

class Armazenamento:
    def salvar(self, pedido: Pedido): ...

class Enviador:
    def enviar_email(self, pedido: Pedido, dest: str): ...
\end{lstlisting}

\begin{exampleblock}{Insight}
Cada classe tem \textbf{uma única razão para mudar}. Testar cada uma é simples e independente.
\end{exampleblock}
\end{frame}

\begin{frame}[fragile]{OCP --- O Problema}
\begin{lstlisting}
class CalculadoraDesconto:
    def calcular(self, pedido, tipo):
        if tipo == "percentual":
            return pedido.total() * 0.10
        elif tipo == "fixo":
            return 5.00
        elif tipo == "fidelidade":
            return pedido.total() * 0.15
        # ... cada novo tipo exige elif
\end{lstlisting}

\begin{alertblock}{Violação OCP}
Para adicionar ``combo'', precisa \textbf{modificar} \texttt{CalculadoraDesconto}. A classe não está fechada para modificação.
\end{alertblock}
\end{frame}

\begin{frame}[fragile]{OCP --- A Solução: Polimorfismo}
\begin{lstlisting}
from abc import ABC, abstractmethod

class Desconto(ABC):
    @abstractmethod
    def calcular(self, total: float) -> float: ...

class DescontoPercentual(Desconto):
    def __init__(self, taxa: float):
        self.taxa = taxa
    def calcular(self, total: float) -> float:
        return total * self.taxa

class DescontoFixo(Desconto):
    def __init__(self, valor: float):
        self.valor = valor
    def calcular(self, total: float) -> float:
        return min(self.valor, total)
\end{lstlisting}

\begin{exampleblock}{Insight}
Novo desconto = nova classe. \texttt{Pedido} não muda. \textbf{Extensão sem modificação.}
\end{exampleblock}
\end{frame}

\begin{frame}[fragile]{LSP --- O Problema Clássico}
\begin{lstlisting}
class Retangulo:
    def __init__(self, largura, altura):
        self._largura = largura
        self._altura = altura

    def set_largura(self, v):
        self._largura = v

    def area(self):
        return self._largura * self._altura

class Quadrado(Retangulo):
    def set_largura(self, v):
        self._largura = v
        self._altura = v   # surpresa!
\end{lstlisting}

\begin{alertblock}{Violação LSP}
Código que espera \texttt{Retangulo} e chama \texttt{set\_largura(5)} seguido de \texttt{area()} obtém resultado inesperado com \texttt{Quadrado}.
\end{alertblock}
\end{frame}

\begin{frame}{LSP --- Regra de Ouro}
\begin{center}
\Large
\textit{``Se S é subtipo de T, objetos de T podem ser\\
substituídos por objetos de S sem alterar\\
as propriedades desejáveis do programa.''}
\end{center}

\vspace{0.5cm}

\textbf{Teste prático:} se o código cliente precisa checar \texttt{isinstance} para decidir o que fazer, provavelmente há violação de LSP.

\begin{exampleblock}{No Podrão}
\texttt{Hamburguer} e \texttt{XTudo} são subtipos de \texttt{Lanche}. Qualquer código que recebe \texttt{Lanche} funciona com ambos --- sem \texttt{if isinstance(...)}.
\end{exampleblock}
\end{frame}

% ================================================================
\section{Padrões de Projeto}
% ================================================================

\begin{frame}[fragile]{Strategy --- O Problema}
\begin{lstlisting}
class Pedido:
    def calcular_desconto(self):
        if self.tipo_desconto == "percentual":
            return self.total * 0.10
        elif self.tipo_desconto == "fixo":
            return 5.00
        elif self.tipo_desconto == "nenhum":
            return 0.0
\end{lstlisting}

Já vimos isso no OCP. A solução: \textbf{Strategy Pattern}.
\end{frame}

\begin{frame}[fragile]{Strategy --- Composição}
\begin{lstlisting}
class SemDesconto(Desconto):
    def calcular(self, total: float) -> float:
        return 0.0

@dataclass
class Pedido:
    cliente: str
    itens: list[str] = field(default_factory=list)
    desconto: Desconto = field(
        default_factory=SemDesconto)

    def total(self, cardapio: dict) -> float:
        bruto = sum(cardapio[i] for i in self.itens)
        return bruto - self.desconto.calcular(bruto)
\end{lstlisting}

\begin{exampleblock}{Composição}
\texttt{Pedido} recebe a estratégia; não sabe qual é. Novo desconto = nova classe, zero mudança em \texttt{Pedido}.
\end{exampleblock}
\end{frame}

\begin{frame}[fragile]{Strategy --- Uso}
\begin{lstlisting}
cardapio = {"X-Tudo": 25.90, "Suco": 8.00}

p1 = Pedido("Ana", desconto=DescontoPercentual(0.10))
p1.adicionar("X-Tudo")
p1.adicionar("Suco")
print(p1.total(cardapio))  # 30.51

p2 = Pedido("Bruno", desconto=DescontoFixo(5.00))
p2.adicionar("X-Tudo")
print(p2.total(cardapio))  # 20.90
\end{lstlisting}

A estratégia pode ser trocada em \textbf{tempo de execução}:
\begin{lstlisting}
p1.desconto = DescontoFixo(10.00)
\end{lstlisting}
\end{frame}

\begin{frame}[fragile]{Strategy Pythônico: Funções como Estratégias}
\begin{lstlisting}
from typing import Callable

@dataclass
class PedidoFunc:
    cliente: str
    itens: list[str] = field(default_factory=list)
    desconto: Callable[[float], float] = lambda t: 0.0

    def total(self, cardapio: dict) -> float:
        bruto = sum(cardapio[i] for i in self.itens)
        return bruto - self.desconto(bruto)
\end{lstlisting}

\begin{lstlisting}
p = PedidoFunc("Carlos",
    desconto=lambda t: t * 0.10)  # 10%
\end{lstlisting}

\begin{exampleblock}{Quando usar}
Classes: quando a estratégia tem estado ou configuração.
Funções: quando é um cálculo simples sem estado.
\end{exampleblock}
\end{frame}

\begin{frame}[fragile]{Factory Method --- O Problema}
\begin{lstlisting}
def criar_lanche(tipo: str):
    if tipo == "xtudo":
        return XTudo()
    elif tipo == "xsalada":
        return XSalada()
    elif tipo == "hamburguer":
        return Hamburguer()
    # ... cresce a cada novo lanche
\end{lstlisting}

\begin{alertblock}{Problemas}
Cadeia de \texttt{if/elif} viola OCP. Adicionar um tipo exige modificar a função.
\end{alertblock}
\end{frame}

\begin{frame}[fragile]{Factory Method --- Registry}
\begin{lstlisting}
class LancheFactory:
    _registry: dict[str, type] = {}

    @classmethod
    def registrar(cls, nome: str, classe: type):
        cls._registry[nome] = classe

    @classmethod
    def criar(cls, nome: str, **kwargs):
        if nome not in cls._registry:
            raise ValueError(f"Tipo desconhecido: {nome}")
        return cls._registry[nome](**kwargs)

# Registro
LancheFactory.registrar("xtudo", XTudo)
LancheFactory.registrar("xsalada", XSalada)
LancheFactory.registrar("hamburguer", Hamburguer)
\end{lstlisting}
\end{frame}

\begin{frame}[fragile]{Factory Method --- Uso e Extensão}
\begin{lstlisting}
# Criação via string (ex: lido de arquivo/API)
lanche = LancheFactory.criar("xtudo", preco=25.90)

# Extensão em runtime --- sem tocar na Factory
class XBacon(Lanche):
    def __init__(self, preco=28.00):
        super().__init__("X-Bacon", preco, ["bacon"])

LancheFactory.registrar("xbacon", XBacon)
lanche2 = LancheFactory.criar("xbacon")
\end{lstlisting}

\begin{exampleblock}{Vantagens}
\begin{itemize}
  \item Desacopla criação de uso
  \item Extensível em runtime (\texttt{registrar})
  \item String $\to$ objeto (útil para configs, APIs, CLIs)
\end{itemize}
\end{exampleblock}
\end{frame}

% ================================================================
\section{Estudo de Caso --- Refatorando o Podrão}
% ================================================================

\begin{frame}[fragile]{O Código ``Legado'' (Procedural)}
\begin{lstlisting}
# Listas paralelas
nomes = ["X-Tudo", "X-Salada", "Suco"]
precos = [25.90, 20.00, 8.00]
quantidades = [0, 0, 0]

def fazer_pedido(idx, qtd):
    quantidades[idx] += qtd

def calcular_total():
    total = 0
    for i in range(len(nomes)):
        total += precos[i] * quantidades[i]
    return total

def aplicar_desconto(total, tipo):
    if tipo == "percentual": return total * 0.9
    elif tipo == "fixo": return total - 5
    return total
\end{lstlisting}
\end{frame}

\begin{frame}{Problemas do Código Legado}
\begin{enumerate}
  \item \textbf{Sopa de variáveis} --- listas paralelas sem conexão explícita
  \item \textbf{SRP violado} --- tudo no mesmo escopo global
  \item \textbf{OCP violado} --- novo desconto $=$ novo \texttt{elif}
  \item \textbf{Sem encapsulamento} --- qualquer parte do código altera \texttt{quantidades}
  \item \textbf{Impossível testar} --- funções dependem de estado global
  \item \textbf{Impossível estender} --- novo lanche $=$ alterar 3 listas
\end{enumerate}

\begin{alertblock}{Diagnóstico}
Este código funciona para 3 lanches. Com 30 lanches, 5 descontos e 3 formas de pagamento, vira pesadelo.
\end{alertblock}
\end{frame}

\begin{frame}[fragile]{Versão OOP Refatorada --- Hierarquia}
\begin{lstlisting}
from abc import ABC, abstractmethod
from dataclasses import dataclass, field

class Lanche(ABC):
    @abstractmethod
    def preco(self) -> float: ...
    @abstractmethod
    def descricao(self) -> str: ...

@dataclass(frozen=True)
class Hamburguer(Lanche):
    nome: str = "Hambúrguer"
    _preco: float = 18.00
    def preco(self) -> float: return self._preco
    def descricao(self) -> str: return self.nome

@dataclass(frozen=True)
class XTudo(Lanche):
    nome: str = "X-Tudo"
    _preco: float = 25.90
    def preco(self) -> float: return self._preco
    def descricao(self) -> str: return self.nome
\end{lstlisting}
\end{frame}

\begin{frame}[fragile]{Versão OOP Refatorada --- Pedido + Strategy}
\begin{lstlisting}
@dataclass
class Pedido:
    cliente: str
    itens: list[Lanche] = field(default_factory=list)
    desconto: Desconto = field(
        default_factory=SemDesconto)

    def adicionar(self, lanche: Lanche):
        self.itens.append(lanche)

    def total(self) -> float:
        bruto = sum(i.preco() for i in self.itens)
        return bruto - self.desconto.calcular(bruto)
\end{lstlisting}

\begin{lstlisting}
class Cozinha:
    def processar(self, pedido: Pedido): ...

class Caixa:
    def finalizar(self, pedido: Pedido): ...
\end{lstlisting}
\end{frame}

\begin{frame}{Mapa de Conceitos: Aulas 1--4}
\begin{table}
\centering
\scriptsize
\begin{tabular}{lll}
\toprule
\textbf{Conceito} & \textbf{Aula} & \textbf{No Podrão} \\
\midrule
Classe e objeto          & 1 & \texttt{Lanche}, \texttt{Pedido} \\
Atributos e métodos      & 1 & \texttt{nome}, \texttt{preco()}, \texttt{adicionar()} \\
Encapsulamento           & 1 & \texttt{\_preco}, \texttt{@property} \\
Herança                  & 2 & \texttt{Hamburguer(Lanche)}, \texttt{XTudo(Lanche)} \\
Classe abstrata          & 2 & \texttt{Lanche(ABC)}, \texttt{Desconto(ABC)} \\
Polimorfismo             & 2 & \texttt{preco()} em cada subtipo \\
Composição               & 3 & \texttt{Pedido} tem lista de \texttt{Lanche} \\
\texttt{@dataclass}      & 4 & \texttt{Pedido}, \texttt{ItemCardapio} \\
\texttt{@classmethod}    & 4 & \texttt{Lanche.from\_csv\_line()} \\
Strategy                 & 4 & \texttt{Desconto} $\to$ \texttt{DescontoPercentual} \\
Factory                  & 4 & \texttt{LancheFactory.criar()} \\
SRP                      & 4 & \texttt{Cozinha}, \texttt{Caixa}, \texttt{Pedido} \\
OCP                      & 4 & Novo desconto $=$ nova classe \\
\bottomrule
\end{tabular}
\end{table}
\end{frame}

% ================================================================
\section{Resumo do Módulo}
% ================================================================

\begin{frame}{As 4 Aulas em Uma Tabela}
\begin{table}
\centering
\scriptsize
\begin{tabular}{clp{7cm}}
\toprule
\textbf{Aula} & \textbf{Tema} & \textbf{Conteúdo Principal} \\
\midrule
1 & Classes e Objetos    & Classe, instância, atributos, métodos, \texttt{\_\_init\_\_}, encapsulamento, \texttt{@property} \\
2 & Herança e Polimorfismo & Herança, \texttt{super()}, ABC, \texttt{@abstractmethod}, \texttt{isinstance}, MRO \\
3 & Composição e Agregação & ``tem-um'' vs ``é-um'', composição, agregação, delegação, \texttt{\_\_contains\_\_} \\
4 & Avançado e SOLID       & \texttt{@dataclass}, \texttt{@classmethod}, \texttt{@staticmethod}, SRP, OCP, LSP, Strategy, Factory \\
\bottomrule
\end{tabular}
\end{table}
\end{frame}

\begin{frame}{Os Três Mantras da POO}
\begin{enumerate}
  \item \textbf{Objetos agrupam dados e comportamento}

  Não use listas paralelas. Crie uma classe que una o dado ao que se faz com ele.

  \item \textbf{Prefira composição a herança}

  Herança cria acoplamento forte. Composição permite trocar comportamento em runtime (Strategy).

  \item \textbf{Cada classe, uma responsabilidade + extensão sem modificação}

  SRP + OCP. Classes pequenas, coesas, extensíveis via polimorfismo.
\end{enumerate}

\begin{exampleblock}{Se lembrar apenas disso\ldots}
Estes três princípios guiam 90\% das decisões de design em POO.
\end{exampleblock}
\end{frame}

\begin{frame}{Próximos Passos}
\begin{columns}[T]
\column{0.48\textwidth}
\textbf{Nesta disciplina:}
\begin{itemize}
  \item Testes unitários (\texttt{pytest})
  \item Type hints e \texttt{mypy}
  \item Tratamento de exceções com classes
  \item Iteradores e geradores
\end{itemize}

\column{0.48\textwidth}
\textbf{Em disciplinas futuras:}
\begin{itemize}
  \item Interface Segregation (I)
  \item Dependency Inversion (D)
  \item Mais padrões: Observer, Decorator, Adapter
  \item Arquitetura de software
\end{itemize}
\end{columns}

\vspace{0.5cm}
\begin{center}
\Large \textbf{Obrigado! Dúvidas?}

\normalsize \texttt{andreyoshiaki@dcomp.ufs.br}
\end{center}
\end{frame}

\end{document}
